\chapter{Construire un flux de travail de recherche}\label{ch:rs:build-a-research-workflow}

Un résultat vieux de dix-huit mois ne peut pas être reproduit. Le fournisseur a depuis révisé son historique, le code a évolué, et personne n'a conservé la graine aléatoire du
bootstrap qui a produit l'intervalle de confiance figurant dans les diapositives du comité. Chacun des trois éléments est banal, et ensemble ils rendent le résultat inconnaissable~: il peut
avoir été juste, mais personne ne peut le dire. Ce dernier chapitre du livre reprend son étude de référence, depuis un \omterm{def:rs:build-a-research-workflow:repro}{instantané de données} jusqu'à la fiche de synthèse, dans un pipeline dont chaque
étape est indexée par le hachage de son code, de ses paramètres, de sa graine et de ses entrées, si bien qu'un résultat est un manifeste qu'on peut réexécuter à l'identique, et qu'une modification, où qu'elle soit, dit
exactement ce qu'elle a changé. La construction est \texttt{firm.workflow}.

\section{Reproductibilité}

\begin{definition}[Résultat reproductible, instantané de données, stockage adressé par contenu, étape du pipeline, verrouillage de l'environnement]\label{def:rs:build-a-research-workflow:repro}
Un \emph{résultat reproductible}\index{résultat reproductible} est celui que les mêmes données, le même code et les mêmes réglages reproduisent à l'identique, bit pour bit (livre~4, chapitre~25), une norme minimale
lorsqu'une réplication indépendante n'est pas possible (Peng). Un \emph{instantané de données}\index{instantané de données} est une copie immuable des données telles qu'elles ont été utilisées, identifiée
par son contenu. Le \emph{stockage adressé par contenu}\index{stockage adressé par contenu} range chaque objet sous le hachage de son contenu, de sorte qu'un même contenu ait
toujours la même adresse, et que deux contenus différents ne puissent pas partager la même (le magasin d'objets de Git en est l'exemple familier). Une \emph{étape du pipeline}\index{etape du pipeline@étape du pipeline} est une fonction de ses
entrées, de ses paramètres et de sa graine, sans autre source de variation. Un \emph{verrouillage de l'environnement}\index{verrouillage de l'environnement} consigne les versions exactes de l'interpréteur
et des bibliothèques avec lesquelles un résultat a été produit.
\end{definition}

Sandve et ses coauteurs ont ramené la pratique à dix règles~; la première, la sixième et la neuvième en portent l'essentiel~: pour chaque résultat, consigner la manière dont il a été produit~; pour
les analyses qui comportent du hasard, noter les graines~; relier les énoncés du texte aux résultats qui les fondent. Le livre applique la neuvième règle à lui-même depuis
le chapitre~1~: chaque nombre imprimé dans ces chapitres est vérifié par un test qui le recalcule. Ce qu'un desk de recherche ajoute, c'est l'échelle~: de nombreuses études, des étapes partagées,
des données de fournisseurs qui changent sous leurs pieds. La réponse consiste à faire en sorte que ce soit le pipeline, et non l'analyste, qui se souvienne de la manière dont chaque résultat a été produit.

\texttt{firm.workflow} indexe chaque étape par le SHA-256 de son nom, le hachage de son code source, ses paramètres, sa graine et les hachages de contenu des
sorties de ses entrées, stocke la sortie sous cette clé, puis en calcule le hachage (\cref{lst:rs:build-a-research-workflow:run}). Une étape dont la clé est déjà stockée
ne s'exécute pas. Comme une clé dépend des sorties de ses entrées et non de leurs clés, un changement en amont qui laisse un résultat intermédiaire inchangé s'arrête à cet endroit. Une exécution
écrit un manifeste~: pour chaque étape, sa clé, le hachage du code, ses paramètres, sa graine, les hachages des entrées, le hachage de la sortie et l'indication de son exécution effective, ainsi que l'environnement.

\section{Revue}

\begin{definition}[Liste de contrôle de revue]\label{def:rs:build-a-research-workflow:review}
Une \emph{liste de contrôle de revue}\index{liste de contrôle de revue} est la liste des questions auxquelles une \omterm{def:rs:the-research-process:gate}{revue de recherche} (chapitre~1) répond avant qu'un résultat soit accepté~: qu'il soit reproduit à partir de son manifeste,
que ses données soient point-in-time, que ses essais soient comptabilisés, que son \omterm{def:rs:vectorised-backtests:fidelity}{niveau de fidélité} corresponde à l'affirmation, et que ses coûts et
sa \omterm{def:rs:capacity-decay-and-crowding:capacity}{capacité} soient ceux de la taille à laquelle il sera mis en œuvre.
\end{definition}

Les chapitres du livre constituent sa liste de contrôle, et un relecteur peut poser leurs questions à n'importe quel manifeste.
\begin{itemize}
\item \textbf{Reproduit.} Le manifeste se reproduit-il à l'identique dans un cache vide, dans l'environnement verrouillé~? Quelles étapes n'ont pas de graine~?
\item \textbf{Point-in-time.} L'instantané est-il identifié, et ses révisions sont-elles consignées (chapitre~3)~? L'univers et les expositions sont-ils connus à chaque date
(chapitres~4 et~24)~?
\item \textbf{Comptabilisé.} Combien d'essais figurent dans la famille du \omterm{def:rs:the-research-process:log}{journal de recherche}, et quelles sont les valeurs du ratio de Sharpe déflaté et de la \omterm{def:rs:overfitting:pbo}{probabilité de surapprentissage du backtest}
(chapitres~1 et~20)~?
\item \textbf{Fidèle.} À quel \omterm{def:rs:vectorised-backtests:fidelity}{niveau de fidélité} le \omterm{def:rs:vectorised-backtests:backtest}{backtest} a-t-il été réalisé, et l'affirmation reste-t-elle dans ce que ce niveau permet de dire (chapitres~16 à~19)~?
\item \textbf{Coûté et dimensionné.} Les coûts sont-ils ceux du modèle ajusté à la taille visée, et où se situe la \omterm{def:rs:capacity-decay-and-crowding:capacity}{capacité} (chapitres~27 et~28)~?
\item \textbf{Énoncé.} Chaque chiffre du rapport est-il rattachable à la sortie d'une étape, et chaque graphique à ses données (neuvième et septième règles de Sandve)~?
\end{itemize}
Arnott, Harvey et Markowitz ont proposé un protocole comparable pour la recherche en apprentissage automatique~; le pipeline rend mécaniques ses premières questions, si bien que la revue
peut consacrer son temps aux dernières.

\section{Construction~: le pipeline, des données brutes à un résultat révisé}

L'étude de référence sur \texttt{firm.synthmkt} (\cref{fig:rs:build-a-research-workflow:dag}, \cref{lst:rs:build-a-research-workflow:make}) comporte huit étapes~:
l'\omterm{def:rs:build-a-research-workflow:repro}{instantané des données}, avec les révisions éventuelles apportées par le fournisseur~; les 500 noms les plus liquides de chaque jour~; deux caractéristiques, le momentum de 12 à 1 mois et l'inversion à un jour~;
leurs statistiques de \omterm{def:rs:anatomy-of-a-predictor:card}{fiche de prédicteur}~; un mélange à parts égales~; un portefeuille mensuel neutre en dollars~; un \omterm{def:rs:vectorised-backtests:backtest}{backtest} de niveau 1 avec des coûts de dix points de base~; et une fiche de synthèse avec un
bootstrap par blocs du ratio de Sharpe, à graine fixée. La première exécution calcule chaque étape en quelques secondes~; la seconde lit tout depuis le cache.

\begin{omfigure}
\begin{tikzpicture}[st/.style={draw=omDef,fill=omDef!8,rounded corners=2pt,font=\scriptsize,minimum width=1.7cm,minimum height=0.7cm,align=center},
  hit/.style={st,draw=omThm,fill=omThm!12},arr/.style={-{Stealth},thick,black!60}]
\node[hit] (s) at (0,0) {instantané};
\node[hit] (u) at (2.1,0) {univers};
\node[hit] (f) at (4.2,0) {caractéristiques};
\node[hit] (c) at (6.3,1.3) {fiches};
\node[st] (b) at (6.3,0) {mélange};
\node[st] (p) at (8.4,0) {portefeuille};
\node[hit] (l) at (10.5,0) {niveau 1};
\node[st] (t) at (12.6,0) {fiche de synthèse};
\foreach \a/\b in {s/u,u/f,f/c,f/b,b/p,p/l,l/t} \draw[arr] (\a) -- (\b);
\draw[arr] (s.north) to[out=50,in=130] (f.north);
\draw[arr] (s.north east) to[out=70,in=180] (c.west);
\draw[arr] (u.south) to[out=-50,in=-130] (b.south);
\draw[arr] (u.south east) to[out=-60,in=-120] (p.south);
\draw[arr] (s.south) to[out=-40,in=-140,looseness=0.6] (l.south);
\end{tikzpicture}

\omcaption{Les huit étapes de l'étude de référence. Une révision du fournisseur portant sur un nom qui n'entre jamais dans l'univers relance les cinq étapes en rouge, mais seule la sortie de l'instantané
change~: les quatre autres reproduisent leurs sorties précédentes, et le mélange, le portefeuille et la fiche de synthèse sont lus depuis le cache.}\label{fig:rs:build-a-research-workflow:dag}
\end{omfigure}

Le résultat propre de l'étude est modeste, et le pipeline le présente avec son incertitude~: le \omterm{def:rs:anatomy-of-a-predictor:ic}{coefficient d'information de rang} (rank IC) quotidien moyen du momentum vaut 0.0076 ($t = 1.39$), celui de l'inversion
0.0379 ($t = 16.05$) sur 2\,267 jours, et le portefeuille mensuel issu du mélange a un ratio de Sharpe après coûts de 0.08 (erreur type 0.33~; intervalle bootstrap allant de
$-0.67$ à 0.74), avec un renouvellement de 3.5~\% de son capital par jour. Le signal quotidien fort de l'inversion ne survit pas à un portefeuille reconstruit chaque mois, le décalage d'horizon
du chapitre~6. L'essentiel ici est la manière dont le résultat se comporte face au changement~:

\begin{center}
\small
\begin{tabular}{@{}p{5.0cm}p{6.6cm}p{3.6cm}@{}}
\toprule
changement & étapes relancées & sorties modifiées\\
\midrule
aucun & aucune & aucune\\
révision, nom jamais dans l'univers & instantané, univers, caractéristiques, fiches, niveau 1 & instantané\\
révision, nom toujours dans l'univers & les huit & les huit\\
poids du mélange 0.7 et 0.3 & mélange, portefeuille, niveau 1, fiche de synthèse & ces quatre-là\\
le code de la fiche de synthèse gagne un \omterm{def:rs:performance-measurement:hit}{taux de réussite} (hit rate) & fiche de synthèse & fiche de synthèse\\
\bottomrule
\end{tabular}
\end{center}

Une révision portant sur un nom qui n'est jamais parmi les 500 plus liquides relance chaque étape qui lit l'instantané, et chacune d'elles reproduit sa sortie précédente, si bien que rien
en aval ne s'exécute~: la différence des manifestes montre que les paramètres et la sortie de l'instantané ont changé, et que seules les entrées des autres ont changé. Une révision à l'intérieur de l'univers modifie
tout, et le signale. Une modification du code ne change que son étape. Et la reproduction, c'est-à-dire la réexécution de l'ensemble du manifeste dans un cache vide, renvoie le même hachage de sortie
pour les huit étapes~; si l'on omet la graine du bootstrap, elle signale exactement une étape, la fiche de synthèse (\cref{lst:rs:build-a-research-workflow:repro}). Enfin, l'exécution est
enregistrée comme un essai dans le \omterm{def:rs:the-research-process:log}{journal de recherche} du chapitre~1, avec le hachage de son manifeste comme identité de code et le hachage de l'instantané comme identité de données~: le
résultat vieux de dix-huit mois, s'il avait été produit de cette façon, ne serait qu'une ligne du journal et une seule commande.

\section{Tutoriel~: le résultat qui n'a pas pu être reproduit}\label{tut:rs:build-a-research-workflow:tut}

\begin{tutorial}
\textbf{Objectif.} Exécuter l'étude de référence avec \texttt{firm.workflow}, modifier tour à tour les données, les paramètres et le code, reproduire le résultat bit à bit, et
l'enregistrer. \textbf{État final~:} le tableau et la \cref{fig:rs:build-a-research-workflow:dag}.
\begin{tutsteps}
\item \textbf{Exécution d'un pipeline}~: clés issues du code, des paramètres, de la graine et des hachages des entrées~; le cache~; le manifeste.
\omcode{code/firm/workflow/firm_workflow.py}{107}{123}{Une exécution de pipeline adressée par contenu.}\label{lst:rs:build-a-research-workflow:run}
\item \textbf{Reproduction et différence.}
\omcode{code/firm/workflow/firm_workflow.py}{126}{142}{Réexécution à partir de zéro, et ce qui a changé entre deux manifestes.}\label{lst:rs:build-a-research-workflow:repro}
\item \textbf{L'étude de référence}, en huit étapes.
\omcode{code/research/29-build-a-research-workflow/python/rs_workflow.py}{115}{129}{Le pipeline de l'étude.}\label{lst:rs:build-a-research-workflow:make}
\item \textbf{Exécuter} \texttt{rs\_workflow.make()} dans un répertoire de cache~; réviser un nom avec \texttt{never\_in\_universe} et \texttt{always\_in\_universe}~; modifier les
pondérations~; remplacer par \texttt{tearsheet\_v2}~; appeler \texttt{reproduce} et \texttt{register}.
\end{tutsteps}
\textbf{Pour aller plus loin.} Ajouter une étape de niveau 2 (chapitre~17) sur le plus grand nom du livre et observer quels changements l'atteignent~; stocker le cache dans un répertoire partagé et
laisser un collègue reproduire le manifeste sur une autre machine, après avoir comparé d'abord les environnements.
\end{tutorial}

\section{Construction~: le workflow}\label{bld:rs:build-a-research-workflow:workflow}

\begin{build}
\bfield{Rôle} Chaque résultat sur lequel s'appuie la maison est produit par un pipeline dont le manifeste le reproduit, indique ce dont il dépend et est enregistré dans le
\omterm{def:rs:the-research-process:log}{journal de recherche}.
\bfield{Interface} \texttt{content\_hash(obj)}, \texttt{Stage(name, fn, deps, params, seed)}, \texttt{Pipeline(stages, cache\_dir).run(target)} renvoyant les sorties et
un manifeste, \texttt{environment()}, \texttt{reproduce(make\_pipeline, manifest, cache\_dir)}, \texttt{diff(m1, m2)}, \texttt{register(manifest, log, family, ts,
metrics)}.
\bfield{Règles} Les étapes sont des fonctions pures des entrées, des paramètres et de la graine~; les données n'entrent que par des étapes d'instantané dont les révisions sont des paramètres~; toute étape qui
tire des nombres aléatoires prend une graine~; les manifestes et les caches sont conservés avec le résultat~; un résultat se relit à partir de son manifeste, et non à partir de ses diapositives.
\bfield{Tests d'acceptation} \texttt{code/firm/workflow/tests/}~: un hachage canonique (l'ordre des clés est indifférent, les types sont significatifs)~; une seconde exécution entièrement issue du cache~; une coupure
précoce lorsqu'une révision est écartée~; des différences nommant les champs modifiés~; une reproduction bit à bit, et une reproduction échouée signalant exactement l'étape sans graine~;
l'enregistrement dans le \omterm{def:rs:the-research-process:log}{journal de recherche}.
\bfield{Pour aller plus loin} Hachages de code qui suivent les fonctions appelées par une étape~; caches distants partagés par le desk~; \omterm{def:rs:build-a-research-workflow:repro}{verrouillages de l'environnement} imposés avant la reproduction.
\end{build}

\begin{omsources}
\item R.~D.~Peng, «~Reproducible research in computational science~», \emph{Science} 334(6060), 2011.
\item G.~K.~Sandve, A.~Nekrutenko, J.~Taylor et E.~Hovig, «~Ten simple rules for reproducible computational research~», \emph{PLoS Computational Biology} 9(10),
2013.
\item \emph{Pro Git}, 2e édition, section 10.2, «~Git internals: Git objects~» (git-scm.com).
\item R.~Arnott, C.~R.~Harvey et H.~Markowitz, «~A backtesting protocol in the era of machine learning~», \emph{Journal of Financial Data Science} 1(1), 2019.
\end{omsources}

\section{Exercices}

\begin{exercise}[$\star$]\label{exo:rs:build-a-research-workflow:1}
Expliquer pourquoi la clé d'une étape doit inclure le hachage de son code, et ce qui manque à une clé ne contenant que ses paramètres.
\end{exercise}

\begin{exercise}[$\star$]\label{exo:rs:build-a-research-workflow:2}
Expliquer pourquoi les clés sont construites à partir des hachages de sortie des entrées plutôt qu'à partir des clés des entrées.
\end{exercise}

\begin{exercise}[$\star$]\label{exo:rs:build-a-research-workflow:3}
Le fournisseur révise un rendement d'une action qui n'a jamais été dans l'univers. Indiquer quelles étapes s'exécutent et quelles sorties changent.
\end{exercise}

\begin{exercise}[$\star\star$]\label{exo:rs:build-a-research-workflow:4}
Indiquer lesquelles des dix règles de Sandve le pipeline fait respecter, et lesquelles il laisse aux personnes.
\end{exercise}

\begin{exercise}[$\star\star$]\label{exo:rs:build-a-research-workflow:5}
Une étape appelle une fonction auxiliaire située dans un autre module. Son hachage de code est le hachage du seul code source de l'étape. Expliquer ce qui peut mal tourner, et comment le corriger.
\end{exercise}

\begin{exercise}[$\star\star$]\label{exo:rs:build-a-research-workflow:6}
La reproduction échoue sur la machine d'un collègue pour chaque étape qui utilise des réductions en virgule flottante. Indiquer ce qu'il faut vérifier en premier.
\end{exercise}

\begin{exercise}[$\star\star\star$]\label{exo:rs:build-a-research-workflow:7}
\emph{Programmation.} Supprimer la graine de la fiche de synthèse et exécuter \texttt{reproduce}. Puis donner au bootstrap une graine dérivée des autres clés du manifeste, et expliquer pourquoi cette graine est
reproductible tout en restant honnête.
\end{exercise}

\begin{exercise}[$\star\star\star$]\label{exo:rs:build-a-research-workflow:8}
\emph{Trouver la faille.} «~L'étude est reproductible~: le carnet est sous contrôle de version.~»
\end{exercise}

\section{Problème~: le résultat qui n'a pas pu être reproduit}

\begin{problem}[{Problème du week-end --- un résultat conservé}]\label{pb:rs:build-a-research-workflow:1}
L'étude de référence dans \texttt{firm.workflow}.

\medskip\noindent\textbf{Partie I --- Le pipeline.}
\begin{enumerate}
\item Lister les huit étapes et leurs dépendances.
\item Préciser ce qui entre dans la clé d'une étape, et ce qui entre dans le manifeste.
\item Expliquer ce que fait la seconde exécution, et pourquoi.
\item Indiquer ce que conclut l'étude, avec son incertitude.
\end{enumerate}

\medskip\noindent\textbf{Partie II --- Modifications.}
\begin{enumerate}[resume]
\item Indiquer quelles étapes sont relancées après une révision portant sur un nom jamais présent dans l'univers, et quelles sorties changent.
\item Faire de même après une révision portant sur un nom toujours présent dans l'univers.
\item Faire de même après une modification des pondérations de la combinaison.
\item Faire de même après une modification du code de la fiche de synthèse.
\end{enumerate}

\medskip\noindent\textbf{Partie III --- Reproduction.}
\begin{enumerate}[resume]
\item Indiquer ce que renvoie la reproduction du manifeste dans un cache vide.
\item Expliquer ce qui se passe sans la graine du bootstrap.
\item Indiquer ce que consigne l'environnement, et quand cela importe.
\end{enumerate}

\medskip\noindent\textbf{Partie IV --- Le verdict.}
\begin{enumerate}[resume]
\item Énoncer le \emph{résultat nommé}~: les étapes invalidées par une révision du fournisseur, et la différence de manifestes qui prouve la reproduction.
\item Expliquer comment l'exécution est enregistrée dans le \omterm{def:rs:the-research-process:log}{journal de recherche}.
\item Indiquer quels points de la \omterm{def:rs:build-a-research-workflow:review}{liste de contrôle de revue} le pipeline couvre, et lesquels il ne couvre pas.
\item Expliquer pourquoi le signal quotidien du retournement ne survit pas au portefeuille mensuel.
\item Expliquer comment le résultat vieux de dix-huit mois aurait été conservé.
\item Indiquer ce qu'il faudrait ajouter avant que le desk s'appuie sur le pipeline.
\item Définir en une phrase ce qu'est un \omterm{def:rs:build-a-research-workflow:repro}{résultat reproductible}.
\end{enumerate}
\end{problem}

\section{Questions d'entretien}

\begin{interviewq}[$\star$ \iqroles{chercheur, développeur}]\label{iq:rs:build-a-research-workflow:1}
Un \omterm{def:rs:vectorised-backtests:backtest}{backtest} que vous avez lancé l'an dernier donne aujourd'hui un ratio de Sharpe différent. Énumérer les causes possibles.
\end{interviewq}

\begin{interviewq}[$\star\star$ \iqroles{chercheur}]\label{iq:rs:build-a-research-workflow:2}
Quelles pratiques rendent la recherche computationnelle reproductible~?
\end{interviewq}

\begin{interviewq}[$\star\star$ \iqroles{développeur}]\label{iq:rs:build-a-research-workflow:3}
Comment concevriez-vous un cache pour les étapes d'un pipeline de recherche~? Comment l'invalideriez-vous~?
\end{interviewq}

\begin{interviewq}[$\star\star$ \iqroles{chercheur, risque}]\label{iq:rs:build-a-research-workflow:4}
Que vérifieriez-vous lors de la revue de la recherche de stratégie d'un collègue~?
\end{interviewq}

\begin{interviewq}[$\star\star$ \iqroles{développeur}]\label{iq:rs:build-a-research-workflow:5}
Comment rendriez-vous reproductible, d'une machine à l'autre, un calcul qui utilise des nombres aléatoires~?
\end{interviewq}

\begin{interviewq}[$\star\star\star$ \iqroles{développeur, chercheur}]\label{iq:rs:build-a-research-workflow:6}
Votre fournisseur retraite cinq ans d'historique. Comment trouveriez-vous chaque résultat qui en dépendait, et lesquels changent~?
\end{interviewq}
